\subsection{A uniform bound from the excess number of arrows}

\begin{theorem}\label{C:arrow-excess-four}
Let $C$ be a finite category with $k$ objects and $n$ arrows.
If $n\leq2k+3$, then $BC$ is not homotopy equivalent to
$\mathbb CP^2$.
In particular, every category with at least nine objects and
at most twenty-one arrows is excluded.
\end{theorem}
\begin{proof}
Only connected categories need consideration.  Their
undirected support has at least $k-1$ edges, and their
identities supply $k$ further arrows.  Define the nonnegative
excess
\[
 e=n-(2k-1)\leq4.
\]
Delete any object whose identity factors through another
object, using the already proved retract deletion.  This
preserves the classifying-space homotopy type and
connectedness.  At least its identity and two cross arrows
are removed; hence $k$ drops by one and $e$ drops by at
least one.  We may therefore assume that no object is a
retract of another, still with $e\leq4$.
If at most two objects remain, the arrow count is at most
seven, already excluded by Theorem~\ref{C:CP2-lower16}.

We first remove the possibility of a nontrivial strongly
connected component.  Let such a component have $t\geq2$
objects.  Every ordered cross hom-set inside it is nonempty.
Every corner inside it has at least two arrows: otherwise
the identity in that trivial corner factors through any
other object of the component.  Thus its full subcategory
contains at least $t^2+t$ arrows.  If it is the whole category,
then $t=k\geq3$ gives
\[
 e\geq k^2-k+1\geq7,
\]
which is impossible.

Suppose it is proper.  The other objects contribute at least
$k-t$ identities.  Contract the component in the undirected
support.  A connected graph on the remaining $k-t+1$
vertices needs at least $k-t$ edges.  At least one such edge
meets the contracted component, and any arrow across that
interface forces an arrow to or from every one of its $t$
objects, by composition inside the component.  Consequently
the arrows not internal to the component include at least
$k-1$ cross arrows.  In total
\[
 n\geq(t^2+t)+(k-t)+(k-1)=2k+t^2-1,
 \qquad e\geq t^2.
\]
Equality is therefore forced, with $t=2$ and $e=4$.
All inequalities used must be equalities.  There is exactly
one neighboring object $v$ outside the component; it has a
trivial corner, and the two cross hom-sets at this interface
are singletons, both oriented toward $v$ or both away from
$v$.  All other corners are trivial, and the contracted
support is a tree with singleton cross hom-sets.

The full subcategory on the component and $v$ has $v$ as a
terminal or initial object.  It meets the rest of the
category only in $v$; a nerve simplex cannot use an object
of the component and any other outside object.  Its
contractible nerve can therefore be removed by a homotopy
pushout along that point.  The remaining nerve is the tree
of singleton arrows on trivial corners, hence contractible.
This case cannot be a projective plane.

We are left with acyclic support.  Put
\[
 D=\sum_x(|\operatorname{End}(x)|-1),\qquad
 P=\sum_{x\ne y,\ C(x,y)\ne\varnothing}(|C(x,y)|-1),
 \qquad R=E-k+1,
\]
where $E$ is the number of edges of the undirected support.
There is at most one direction on each such edge, and
\begin{equation}\label{C:excess-decomposition}
 e=D+P+R\leq4.
\end{equation}
A strict chain of four arrows would supply all ten edges
on its five objects, and hence force $R\geq6$.
Thus there are at most three strict changes of object in
any simplex.

If all corner algebras are semisimple over $\mathbb Q$,
the relative bar over their product has no internal higher
Tor.  Its strict-object degree is at most three, and
$H_4(BC;\mathbb Q)=0$.

Every corner has at most five elements.  If one has five,
then $D=4$ and $P=R=0$.  There is only that one nontrivial
corner, the support is a tree, and every cross hom-set is
a singleton.  A valid acyclic category on a tree has no
composable pair of strict arrows, since their composite
would create a triangle.  The five-element corner is
incident to a singleton leg to a point.  The corresponding
two-object classifying space is the join of its corner
space with a point, so every global degree-two class
restricts to zero on that corner.  It also restricts to
zero on every other vertex.  The quotient by the union
of vertex spaces is a wedge of suspended edge bars, as
in Corollary~\ref{C:height-one-zero-H2-products}.
All the degree-two classes are relative and all their
products vanish.

Henceforth every corner has at most four elements.  By
Theorem~\ref{E:all-six-types} its classifying space is
rationally acyclic.  If $R=0$, support is again a tree of
height one, and Corollary~\ref{C:height-one-zero-H2-products}
applies.  If $P=0$, collapse every corner to its object.
This actual functor to the support poset is a rational
homology equivalence.  To see the comparison directly,
filter by the finite number of strict object changes:
because every cross hom-set is a singleton, each
associated-graded correspondence bar is the tensor
product of the ordinary bars of the corners in that
chain.  All those corner bars are rationally acyclic.
The comparison is therefore an isomorphism on the first
page and on total homology.  The support poset has
dimension at most three, so its fourth homology vanishes.

The only case still requiring consideration has a
nonsemisimple corner algebra and $P,R\geq1$.
An order-one or order-two monoid algebra is semisimple,
so $D\geq2$.  Equation~\eqref{C:excess-decomposition}
forces
\[
 D=2,\qquad P=1,\qquad R=1.
\]
There is exactly one corner of order three, all others
are trivial, the support is unicyclic, and exactly one
cross hom-set has two arrows; all other nonempty cross
hom-sets are singletons.

If the unique undirected cycle is not a triangle,
there can be no composable pair of strict arrows, and
the height-one corollary again applies.  Otherwise the
unique cycle is a transitive triangle.  The category
nerve is the union of its full triangle subcategory
and its full two-object bridge subcategories.  Any
simplex with two strict changes of object must lie in
that triangle; every other simplex lies in a bridge
piece or a corner.  These pieces form a tree of
subcomplexes meeting in their corner nerves.  The
intersection spaces are rationally acyclic.  The
Mayer--Vietoris sequence thus identifies fourth rational
homology of the whole nerve with the direct sum of
fourth homology of these pieces.

Each bridge piece has trivial corners except possibly
the single order-three corner.  Both trivial modules
of an order-three monoid algebra have projective
dimension at most one, as established in the proof
of Proposition~\ref{C:semisimple-ends-small-middle}.
Its edge bar has no homology in internal degree at
least two, and the bridge collage has zero fourth
homology.

If the triangle has only trivial corners, its nerve
has dimension at most two.  If its order-three corner
is an endpoint, the strict-object bar has internal
degree at most one, again by the trivial-module
projective-dimension bound at that endpoint; all other
corner algebras in the piece are $\mathbb Q$.
Total degree four is absent.  Finally, if the
order-three corner is in the middle, at least one
adjacent hom-set is a singleton, since there is only
one doubled hom-set in the entire category.  The
singleton-cofinal deletion removes that middle
object, leaving a graph.  Thus the triangle piece
also has zero fourth rational homology.
The direct-sum comparison proves $H_4(BC;\mathbb Q)=0$
in the last case.  Every alternative has now excluded
a projective-plane square or its fourth homology,
which proves the theorem.
\end{proof}
